%==================================================
%      PREAMBOLO e DICHIARAZIONI INIZIALI
%==================================================
\documentclass[10pt,oneside,a4paper]{article}

\usepackage[latin1]{inputenc} 
\usepackage[italian]{babel}
\usepackage[T1]{fontenc}
\usepackage{siunitx} %Inserisce automaticamente i dati con le unit�  di misura correttamente formattate del SI (utilizzo: \SI{0.82}{m^2}, in generale \SI{misura con il punto decimale}{unit�  di misura})
\sisetup{output-decimal-marker = {.}, separate-uncertainty = true, input-uncertainty-signs = \pm, detect-weight=true, detect-family=true} %per usare SI con il punto decimale
\usepackage{listings} %Per citare codice informatico formattandolo correttamente
\usepackage{amsmath,amsthm,verbatim,amssymb,amsfonts,amscd,graphicx,mathtools}
\usepackage[makeroom]{cancel}
\newcommand{\abs}[1]{\left\lvert\,#1\,\right\rvert}
\usepackage{geometry}
\usepackage{epigraph}
\usepackage{booktabs}	%tabelle migliorate
\usepackage{tablefootnote}	%note a pi� di pagina in tabella
\usepackage{threeparttable} %tabella con note a pi� di tabella
\usepackage{caption}	%descrizione per figure
\usepackage{dblfnote}
\captionsetup{tableposition=top,figureposition=bottom,font=small} %setup descrizione
\usepackage{float}
\usepackage{esvect} %vettori
\usepackage{longtable} %tabelle lunghe
\usepackage[dvipsnames]{xcolor}
\definecolor{sepia}{HTML}{80002A}
\usepackage[colorlinks=true, citecolor=black, linkcolor=sepia, urlcolor=black]{hyperref}
\usepackage{mathrsfs}
\usepackage{circuitikz}


\usepackage{multicol}
\newenvironment{Figure}
  {\par\medskip\noindent\minipage{\linewidth}}
  {\endminipage\par\medskip}

\newcommand{\var}{\operatorname{var}}
\newcommand{\cov}{\operatorname{cov}}


\usepackage{listings} %Per inserire codice
\lstnewenvironment{codice_c}[1][]
{\lstset{basicstyle=\small\ttfamily, columns=fullflexible,
keywordstyle=\color{red}\bfseries, commentstyle=\color{blue},
language=C, basicstyle=\small,
numbers=left, numberstyle=\tiny,
stepnumber=2, numbersep=5pt, frame=shadowbox,  showstringspaces=false, #1}}{}

%==================================================
%                  PRIMA PAGINA
%==================================================

\title{\textsc{\textbf{Esercitazione 3}: Amplificatore operazionale 1}}
\author{\small{G. Galbato Muscio} \and \small{L. Gravina} \and \small{L. Graziotto}}
\date{\today}

\begin{document}
	\begin{figure}
		\centering
		\includegraphics[scale=0.5, trim={2.8cm 8.9cm 0 9cm}, clip]{logo.png}
	\end{figure}
	\maketitle
	\begin{center} 
		\fbox{{\fontsize{12pt}{8mm}\textsc{Gruppo 11}}} \\
	\end{center}
\hrule
\vspace{0.5cm}
\renewcommand{\abstractname}{Abstract}
\begin{abstract}
Si misurano lo \emph{slew rate} e il prodotto guadagno per banda passante di un amplificatore operazionale LM358. Quindi, lo si utilizza per realizzare un circuito derivatore.
\end{abstract}
\vspace{4cm}
\tableofcontents %Indice
\newpage


\pagebreak
\begin{multicols}{2}
%==================================================
%      		MISURA DELLO SLEW RATE
%==================================================
\section{Misura dello \emph{slew rate}}
Si utilizza nel seguito l'amplificatore operazionale LM358, e lo si alimenta con una differenza di potenziale di $\pm \SI{12}{V}$, connettendo gli ingressi \texttt{V+} e \texttt{GND} alla doppia alimentazione in continua, positiva e negativa. 

Si realizza, per ottimizzare la misura dello \emph{slew rate}, un \emph{emitter follower}, ossia un amplificatore non invertente con amplificazione unitaria, illustrato nel circuito seguente.

\begin{center}
\begin{circuitikz}
\draw (0,0) node[op amp, yscale=-1] (opamp) {}
(opamp.+) to[short, -*] ++(-0.75,0) node[left] {$V_\text{i}$}
(opamp.-) to[short] ++(0, -1.25) -| (opamp.out)
(opamp.out) to[short, *-*] ++(0.5,0) node[right] {$V_\text{o}$}
;\end{circuitikz}
\end{center}

Si invia all'ingresso non invertente dell'op-amp un segnale rettangolare (onda quadra) di frequenza \SI{111 \pm 111}{Hz} e ampiezza variabile, oscillante tra $0$ e $V_i$. Osservando sull'oscilloscopio il segnale in uscita, si osserva che l'amplificazione � effettivamente circa unitaria (si ha $A_v = \SI{111 \pm 111}{}$, data dal valor medio), ma la risposta presenta un tempo di salita, ovvero non segue istantaneamente l'ingresso. Selezionando la regione crescente lineare e misurandone la pendenza mediante i cursori, si ricava in ciascun caso la \emph{slew rate} come
\[
S = \frac{\Delta V}{\Delta t} = \frac{V_o}{\Delta t}.
\]
Si riportano in tabella~\ref{tab:slewrate} i valori di $V_i$, di $\Delta V$, di $\Delta t$ e della stima di $S$.

\begin{center}
\captionof{table}{Valori misurati per la stima dello \emph{slew rate}}
\label{tab:slewrate}
\begin{tabular}{c|c|c|c}
 $V_i$ [\SI{}{V}] & $\Delta V$ [\SI{}{V}] & $\Delta t$ [\SI{}{s}] & $S$ [\SI{}{V/s}] \\
\hline
\hline
\end{tabular}
\end{center}

La stima dello \emph{slew rate}, data dalla media dei valori misurati con errore associato la deviazione standard, risulta essere
\[
S = \SI{111 \pm 111}{V / s};
\]
tale valore risulta essere compatibile con quello dichiarato dal costruttore, di \SI{0.5}{V / \micro s}. A titolo esemplificativo, si riporta in figura~\ref{fig:slewrate} un'istantanea dell'oscilloscopio, dalla quale risulta visibile il tempo finito impiegato dalla risposta a raggiungere il valore a regime.

%\begin{figure}[H]
%	\begin{center}
%	\includegraphics[width=\linewidth]{}
%	\caption{Stima dello \emph{slew rate} utilizzando un'onda quadra}
%	\label{fig:slewrate}
%	\end{center}
%\end{figure}

Si ripete quindi la misura inviando all'ingresso non invertente un segnale sinusoidale del tipo $V_i = V_p \sin(\omega t)$, di frequenza \SI{111 \pm 111}{Hz} e ampiezza picco-picco $V_p = \SI{111 \pm 111}{V}$. Si stima anche in questo caso lo \emph{slew rate} 
%COMPLETARE: BISOGNA OSSERVARE IL GRAFICO


Poich� la banda passante � in questo caso di \SI{1}{MHz}, come dichiarato dal costruttore, e l'amplificazione � unitaria, si ha che � verificata la richiesta di lavorare sotto la frequenza di taglio, e di rispettare la disuguaglianza
\[
A_v V_p \omega \leq S.
\]  

%==================================================
%    MISURA DEL PRODOTTO GUADAGNO-BANDA PASSANTE
%==================================================
\section{Misura del prodotto guadagno per banda passante}
Si costruisce un amplificatore invertente con medesima alimentazione della sezione precedente ($\pm$\SI{12}{V}) e se ne misura la risposta in funzione della frequenza, per vari valori dell'amplificazione, ottenuti cambiando le resistenze della rete di reazione negativa, ricordando che in questo caso
\[
A_v = -\frac{R'}{R}.
\]
Il circuito utilizzato � il seguente.

\begin{center}
\begin{circuitikz}
\draw (0,0) node[op amp] (opamp) {}
(opamp.+) -| ++(-0.5,-0.5) node[ground]{}
(opamp.out) to[short, *-*] ++(0.5,0) node[right] {$V_\text{o}$}
(opamp.-) to[short] ++(-0.25,0) to[R=$R$] ++(-1.5,0) node[] (Vs) {} ++(0,-2) node[] (ground) {}
(ground) node[ground] {} to[sV = $V_s$, invert] (Vs)
(opamp.-) -| ++(0,1) to[R = $R'$] ++(1.5,0) -| (opamp.out)
(opamp.-) node[circ] {}
;\end{circuitikz}
\end{center}

Le resistenze vengono scelte tenendo conto che si vuole minimizzare la corrente che scorre all'interno dell'OP-AMP (poich� non si ha a che fare con il caso ideale), e che si vuole ottenere una resistenza di ingresso di valore elevato, in modo da avvicinarsi al caso ideale. 

Si invia in ingresso al circuito un segnale sinusoidale di ampiezza picco-picco $V_s$, in modo che sia rispettata la disuguaglianza
\[
V_s \leq \frac{S}{2\pi f A} \simeq \SI{80}{mV},
\]
utilizzando i dati forniti dal costruttore, al fine di evitare un segnale distorto in $V_o$. Sui canali \texttt{CH1} e \texttt{CH2} dell'oscilloscopio vengono verificati rispettivamente il segnale di ingresso e di uscita, e si verifica che quest'ultimo � sfasato di circa $\pi$ rispetto al primo, come previsto dalla dinamica dell'amplificatore invertente.

Si riportano nella tabella~\ref{tab:gainbandwidth} in appendice i valori di $V_s$, $V_o$, dell'amplificazione $A_v = V_o / V_s$ e della frequenza a cui viene effettuata la misura $f$ per i diversi valori delle resistenze $R$ e $R'$ con le quali viene realizzata la rete di reazione. Tali resistenze sono misurate con il ponte, mentre le tensioni e le frequenze con l'oscilloscopio. In figura~\ref{fig:bodeGainbandwidth} � riportato il diagramma di Bode con i punti sperimentali; da questo si ricava il guadagno come media dei valori dell'amplificazione nella regione dove essa � costante, con associata l'incertezza data dall'oscilloscopio (dell'ordine del $3\%$), che risulta dominante rispetto a quella statistica, mentre la banda passante � stimata dall'ascissa del punto di intersezione tra la retta che rappresenta la discesa dell'amplificazione dopo la frequenza di taglio (ottenuta interpolando linearmente i dati) con la retta orizzontale relativa ad un'amplificazione di \SI{-3}{dB}, che rappresenta appunto il valore di $A_v$ alla frequenza di taglio. Si verifica inoltre che la pendenza della retta decrescente sia compatibile con il valore di \SI{-20}{dB/dec}.

%\begin{figure}[H]
%	\begin{center}
%	\includegraphics[width=\linewidth]{}
%	\caption{Diagramma di Bode per la misura del prodotto guadagno per banda passante}
%	\label{fig:bodeGainbandwidth}
%	\end{center}
%\end{figure}

Per le differenti configurazioni della rete di reazione, si ottengono i valori riportati in tabella~\ref{tab:riepilogoGainbandwidth}.

\begin{center}
\begin{table*}
\captionof{table}{Prodotto guadagno per banda passante per le diverse reti di reazione}
\label{tab:riepilogoGainbandwidth}
\begin{tabular}{c|c|c|c|c}
 $R$ [\SI{}{k\ohm}] & $R'$ [\SI{}{k\ohm}] & Guadagno ($A_v$) & Banda passante [\SI{}{kHz}] & Guadagno $\times$ Banda passante [\SI{}{kHz}] \\
\hline
\hline
\end{tabular}
\end{table*}
\end{center}

Si ottiene una stima conclusiva del prodotto guadagno per banda passante, data dalla media pesata delle differenti misurazioni, di 
\[
\text{GBW product} = \SI{111 \pm 111}{kHz},
\]
compatibile con il valore atteso di \SI{1}{MHz} riportato nel datasheet.

%==================================================
%    			CIRCUITO DERIVATORE
%==================================================
\section{Circuito derivatore}
Si realizza un circuito derivatore con opportuna scelta delle componenti $R_1$ e $C$ affinch� si abbia un $\tau = R_1 C$ dell'ordine di $10^{-4}$ \SI{}{s}. Si sceglie la resistenza $R_2$, inoltre, in modo tale che sia $R_2 C \ll R_1 C$. Il circuito utilizzato � il seguente.

\begin{center}
\begin{circuitikz}
\draw (0,0) node[op amp] (opamp) {}
(opamp.+) -| ++(-0.5,-0.5) node[ground]{}
(opamp.out) to[short, *-*] ++(0.5,0) node[right] {$V_\text{o}$}
(opamp.-) to[short] ++(-0.25,0) to[C=$C$] ++(-1,0) to[R=$R_2$] ++(-1.5,0) node[circ] {} node[left] {$V_s$} 
(opamp.-) -| ++(0,1) to[R = $R_1$] ++(1.5,0) -| (opamp.out)
(opamp.-) node[circ] {}
;\end{circuitikz}
\end{center}

L'OP-AMP � alimentato ancora con una tensione di $\pm$ \SI{12}{V}; i valori delle resistenze e della capacit� sono
\[
\begin{aligned}
R_1 &= \SI{111 \pm 111}{k\ohm} \\
R_2 &= \SI{111 \pm 111}{k\ohm} \\
C &= \SI{111 \pm 111}{F},
\end{aligned}
\]
misurati con il multimetro e con il ponte. Si fornisce all'ingresso invertente dell'amplificatore un'onda quadra di ampiezza $V_s$ e frequenza $f$, e si osserva al \texttt{CH2} dell'oscilloscopio una $V_o$ impulsiva, localizzata dove si ha il passaggio dell'onda dalla tensione massima a quella minima; si riporta in figura~\ref{fig:derivatoreOndaquadra} un'istantanea dell'oscilloscopio per questa configurazione.

%\begin{figure}[H]
%	\begin{center}
%	\includegraphics[width=\linewidth]{}
%	\caption{Screenshot dell'oscilloscopio per l'onda quadra in ingresso al derivatore}
%	\label{fig:derivatoreOndaquadra}
%	\end{center}
%\end{figure}

Si studia quindi la risposta in frequenza, inviando un segnale sinusoidale di ampiezza picco-picco $V_s$ e frequenza $f$, e misurando con l'oscilloscopio l'ampiezza $V_o$, quindi si ricava l'amplificazione $A_v = V_o / V_s$. Si riportano nella tabella~\ref{tab:derivatore} in appendice le diverse misure, e nella figura~\ref{fig:bodeDerivatore} il diagramma di Bode dei dati sperimentali. Si osserva, come previsto, che il comportamento del derivatore reale si discosta da quello ideale a causa del taglio dell'amplificazione alle alte frequenze. Si stima la discesa oltre la frequenza di taglio mediante interpolazione lineare, ottenendo una pendenza di \SI{111 \pm 111}{dB/dec}, in accordo con il valore teorico di \SI{-20}{dB/dec}. Per la regione crescente si stima una pendenza di \SI{111 \pm 111}{dB/dec}; poich� in questa si ha che il derivatore realizzato con l'OP-AMP segue l'andamento ideale
\[
\abs{A_v} = \frac{\abs{V_o}}{\abs{V_s}} = R_1 C \omega, 
\]
si confronta tale pendenza con il valore previsto di \SI{20}{dB/dec} e l'intercetta con l'asse delle ascisse, che dal diagramma di Bode � in corrispondenza di $f_0 = \SI{111 \pm 111}{Hz}$, con il valore previsto di $f_0 = 1/ (2\pi R_1 C)$.
%\begin{figure}[H]
%	\begin{center}
%	\includegraphics[width=\linewidth]{}
%	\caption{Screenshot dell'oscilloscopio per l'onda quadra in ingresso al derivatore}
%	\label{fig:bodeDerivatore}
%	\end{center}
%\end{figure}





\end{multicols}
%==================================================
%      				APPENDICE	
%==================================================
\section{Appendice}

\begin{center}
\begin{table}[h]
\captionof{table}{Misure per la stima del prodotto guadagno per banda passante}
\label{tab:gainbandwidth}
\begin{tabular}{c|c|c|c}
$f$ [\SI{}{kHz}] & $V_s$ [\SI{}{mV}] & $V_o$ [\SI{}{V}] & $A_v$ \\
\hline
\hline
\end{tabular}
\end{table}
\end{center}

\begin{center}
\begin{table}[h]
\captionof{table}{Misure per il circuito derivatore}
\label{tab:derivatore}
\begin{tabular}{c|c|c|c}
$f$ [\SI{}{kHz}] & $V_s$ [\SI{}{mV}] & $V_o$ [\SI{}{V}] & $A_v$ \\
\hline
\hline
\end{tabular}
\end{table}
\end{center}





\end{document}